\begin{afterword}
\summaryhead

The author suggests that people who set out to be rationalists are probably more annoyed than average by flaws in reasoning, and that if they want to get anything done together they must tolerate other people's tolerance.

The example is Ben Goertzel, who has something nice to say about everyone, even about someone the post calls ``the most legendary of all AI crackpots.'' The author has come to see this as one of Goertzel's strengths, and argues that to mark him down for it is to insist that he dislike everyone the author dislikes. Cooperation needs some punishment of defectors. Punishing people for their own mistakes is one thing, but punishing those who refuse to punish can lock in an equilibrium harmful to everyone. The author resolves to tolerate people more tolerant than the author, and to judge them only for their own mistakes. A group that demands you hate the right enemies may be the wrong group; but do not demand that your friends be equally hostile to it.

\responsehead

The post's central warning is sound and has support. Punishing defectors can sustain cooperation, but punishing those who will not punish can hold any rule in place, useful or not. Boyd and Richerson's model of 1992 makes this point: strategies that punish non-punishers can make ``any individually costly behavior'' stable, ``whether or not it creates a group benefit.'' The post states this in a sentence and turns it into a personal rule.

The rule is offered as the author's own discipline. The premise, that would-be rationalists are more easily annoyed, is hedged with ``probably'' and qualified at once. The examples are from the author's circle. The post claims no more than a resolution and a warning.

The FAQ the post links for its main example adds that Goertzel later conceded critics' points about the ideas he had praised. The notes give Popper's version of the ``intolerant of intolerance'' view the post dismisses in an aside.

In short: a sound warning against punishing those who will not punish, supported by a model of cooperation and offered as the author's own discipline.
\end{afterword}
